
\subsection{Grupės $\pinp$ ir $\pinm$}

\begin{definition}
	$\pinpm$ yra grupės genruojamos vienetiniu vektorių $\vb v \in R^3$,
	$(\vb v,\vb v) = 1$, juos dauginat algebroje $\Clpm$.
	Gupės elemnetai yra vienanariai $\vb v_1 \vb v_2 \dots \vb v_n$
	ir gupės daugyba sutampa su algebros daugyba.
\end{definition}

\begin{definition}
	Grupė $\spin$ yra pogrupis $\pinp$,
	sudarytas iš visų lyginių $\pinp$ elementų.
\end{definition}

\begin{preposition}
	Jie ir tik jei $v$ priklauso $\pinp$,
	$vv^t=\mathbb1$.
	Jei ir tik jei $v$ priklauso $\pinm$,
	$v\alpha(v^t)=\mathbb1$.
\end{preposition}

\begin{example}
	Patikrinkime, kad $\pinm$ tikrai tenkina grupės aksiomas.
	\begin{itemize}
		\item Asociatyvumas yra paveldimas iš algebros
			$\Clm$,
			kuri yra asociatyvi;
		\item tapatybės elementas yra $\mathbb 1$;
		\item atvikštinis elementas yra $v^{-1}= \alpha(v^t)$;
		\item sandaugos uždarumą patikriname skaičiavimu
	\end{itemize}
	\begin{align*}
		\alpha((vw)^t) vw &= \alpha(w^t v^t) vw && \text{antihomomorfizmo savybė}
		\\
		&= \alpha(w^t) \alpha(v^t) vw && \text{homomorfizmo savybė}
		\\
		&= \alpha(w^t) \mathbb 1 w && \text{iš $\pinm$ apibrėžimo}
		\\
		&= \mathbb 1
		\,.
	\end{align*}
\end{example}

Grupėms $\pinpm$ prikkauso visi vienetiniai vektoriai
perkelti iš $\set R^3$ į $\mathrm{Cl}^\pm(3)$.
Jie yra $\pinpm$ generatoriai,
tad visus $\pinpm$ elementus galima realizuoti,
kaip vienetinių vectorių sandaugą algebroje
$\mathrm{Cl}^\pm(3)$. $\set Z_2$ gradavimas išlika
ir ant $\pinpm$.
Elementai $\vb v$ ir $\vb{vwu}$ yra nelyginiai,
o elementai $\mathbb 1$ ir $\vb{vw}$ yra lyginiai.
Dauginant tarpusavyje lyginius elementus,
gausime tik lyginius elementus,
todėl jie sudaro $\pinpm$ pogrupį.


\begin{preposition}
	Grupių $\pinp$ ir $\pinm$ lyginiai pogrupiai
	yra izomorphiški, kaip grupės.
\end{preposition}
\begin{proof}
	$\pinpm$ lyginiai porgupiai kaip aibės prikauso 
	atitinkamiems $\Clpm$ lyginiams poalgiabriams $A_0^\pm$.
	Poalgebriai $A_0^+ \cong A_0^-$ yra izomorfiški kaip algebros (teiginys \ref{prep:even_iso}),
	tegul $f: A_0^+ \to A_0^-$ žymi izomorphizmą tarp jų.
	Jei $\Clp$ lyginis elementams $v$ priklauso $\pinp$,
	tai $f(v)$ priklauso $\pinm$
	\begin{equation}
		\begin{aligned}
			v^t *_+ v &= \mathbb 1_+
			\\
			f(v^t *_+ v ) &= f\qty(\mathbb 1_+)
			\\
			f( v^t) *_- f(v) &= \mathbb 1_-
			&& \text{izomorfizmo savybė}\\
			{f(v)}^t *_- f(v) &= \mathbb 1_-
			&& \text{transponavimas gerbia isomorfizmą}\\
			\alpha\qty({f(v)}^t) *_- f(v) &= \mathbb 1_-
			&& f \text{ kodomenas yra lyginiai elementai}
			\,,
		\end{aligned}
	\end{equation}
	čia $*_+$ ir $*_-$ yra algebros daugybos atitinkamose algebrose,
	o $\mathbb 1_+$ ir $\mathbb 1_-$ yra 
	atitinkamų algebrų vienetiniai elementai.
	Atitinkamai gauname rezultatą, kad jei $v$ prikauso $\pinp$,
	tai $f^{-1}(v)$ priklauso $\pinm$.
	Apribojus domeną ir kodomeną funksija $f$ yra
	grupių izomorfizmas, nes grupės sandauga
	sutampa su algebros sandauga.
\end{proof}

Tiek $\pinp$, tiek $\pinm$ lyginius pogrupius vadinsime vienodai - $\spin$,
nes jie yra izomofiški.

\begin{preposition}
	Visi $\spin$ elementai yra formos
	\begin{equation}
		a_0 \mathbb1 
		+ a_1 \vb e_{23}
		+ a_2 \vb e_{31}
		+ a_3 \vb e_{12}
	\end{equation}
	kur $a_i \in \set R$ ir
	\begin{equation}
		{a_0}^2
		+{a_1}^2
		+{a_2}^2
		+{a_3}^2
		= 1
	\end{equation}
\end{preposition}
\begin{proof}
	Bendra lyginio $\Clp$ vektoriaus forma yra
	\begin{equation}
		v = 
		  a_0 \mathbb1 
		+ a_1 \vb e_{23}
		+ a_2 \vb e_{31}
		+ a_3 \vb e_{12}
	\end{equation}
	Kad elementas $v$ priklausytų $\spin$, jis turi tenkinti
	$vv^t = \mathbb1$.
	Įsisteatę $v$ bendra formą gauname
	\begin{equation}
		\begin{aligned}
			(a_0 \mathbb1 
			+ a_1 \vb e_{23}
			+ a_2 \vb e_{31}
			+ a_3 \vb e_{12})
			(a_0 \mathbb1 
			- a_1 \vb e_{23}
			- a_2 \vb e_{31}
			- a_3 \vb e_{12})
			&= \mathbb1
			\\
			({a_0}^2
			+{a_1}^2
			+{a_2}^2
			+{a_3}^2)
			\mathbb1
			&= \mathbb1
		\end{aligned}
	\end{equation}
	Rezultatas seka iš fakto kad ${\vb e_{ij}}^2 = -\mathbb1$ 
	ir $\vb e_{ij} \vb e_{kl} = - \vb e_{kl} \vb e_{ij}$.
\end{proof}



\subsection{$\spin$ ir $\pinpm$ veikimas ant $\vecSpace R 3$}

Visos aptartos grupės veikia vektorius iš $\set R^3$, juos pirma patalpinant į
$\mathup{Cl^\pm(V)}$ ir tada veikiant su algebros sandauga
\begin{equation}
	\rho(v).\vb w := \alpha(v) \vb w v^{-1} 
	= \begin{cases}
		\alpha(v) \vb w v^t		& v \in \pinp \\
		v \vb w v^t	& v \in \pinm 
	\end{cases}
\end{equation}

Iškart galima pamatyti kad $\pinpm$ elementai $\vb v$ ir $-\vb v$ veikia vienodai,
nes daugiklis $-1$ du kart iškeliamas prieš algebros sandaugą.
Teiginiai \ref{prep:R_action+} ir \ref{prep:R_action-} parodo kad
vienetinis vektorius $\vb n \in \set R^3$ patalpinti į $\pinpm$,
veikia ant $\vb w \in \set R^3$,
kaip atspindys per plokštumą statmeną vektoriui $\vb n$.
Tada lyginiai elementai $\vb v \vb w$ ant $\set R^3$ veikia,
kaip atspindys per vieną plokštumą tada atspindys per kitą
plokštumą. Tai yra teisingas pasukimas.
Kitas elementas kurio veikimą svarbu patikrinti yra
$\vb e_{123}$, gauname $\rho(\vb e_123).\vb e_i = - e_i$,
tai $e_{123}$ ant $\set R^3$ veikia kaip inversija.





\begin{theorem}[{\cite[Theorem~2.9]{lawson1989spin}}, atvėjis kai $V=\set R^3$ ir $q(v)=\pm(v,v)$]
	Egzistuoja truposios tiksliosios sekos (ang. \emph{short exact sequence})
	\begin{equation}
		\begin{tikzcd}
			1 \arrow[r] &\{1,-1\} \arrow[r]	&\pinp \arrow[r, "\rho"] &\Orth3 \arrow[r] &1
		\end{tikzcd}
	\end{equation}
	\begin{equation}
		\begin{tikzcd}
			1 \arrow[r] &\{1,-1\} \arrow[r]	&\pinm \arrow[r, "\rho"] &\Orth3 \arrow[r] &1
		\end{tikzcd}
	\end{equation}
	\begin{equation}
		\begin{tikzcd}
			1 \arrow[r] &\{1,-1\} \arrow[r]	&\spin \arrow[r, "\rho"] &\SO3 \arrow[r] &1
		\end{tikzcd}
	\end{equation}
	Čia kiekviena rodyklė yra grupių homorfizmas ir kiekvieno homomorfizmo vaizdas
	yra branduolys homomorfizmo einančio po jo sekoje.
\end{theorem}




\subsection{$\spin$ ir $\pinpm$ matriciniai įvaizdžiai}

Dėl patogumo dirbant su grupėmis ir algebromis apibrėšime
nagrinėjamų grupių ir algebrų ištikimus matricinius įvaizdžius.
Šiuos matricinius įvaizdžius užrašysime su Pauli matricų 
ir vienetinės $2\times2$ matricos pagalba
\begin{equation}
	I := \mqty(\imat{2})
	\,,\qquad
	\sigma_1 := \mqty(\pmat{1}) 
	\,,\qquad
	\sigma_2 := \mqty(\pmat{2}) 
	\,,\qquad
	\sigma_3 := \mqty(\pmat{3}) 
	\,.
\end{equation}
Šios matricos tenkina
\begin{align}
	\sigma_i \sigma_i &= I \,,
	&&i\in{1,2,3}\,;
	\\
	\sigma_i \sigma_j &= - \sigma_j \sigma_i\,, 
	&&i\neq j \,;
	\\
	\sigma_i \sigma_j &= i \sigma_k \,,
	&&\text{$i,j,k$ yra indeksų $1,2,3$ cikliniai perstatymai}.
\end{align}


\begin{preposition}
	\label{prep:pauli_rep+}
	Veikimu ant $\Clp$ \emph{lyginio} poalgebrio $A_0^+$ bazės apibrėžtas atvaizdavimas
	\begin{align*}
		\mathbb 1	\mapsto	I		\qquad
		\vb e_{23} 	\mapsto	i\sigma_1	\qquad
		\vb e_{31}	\mapsto	i\sigma_2	\qquad
		\vb e_{12}	\mapsto	i\sigma_3	
	\end{align*}
	yra algebros $\Clp$ \emph{lyginio} poalgebrio ištikimas matricinis įvazdis.
\end{preposition}
\begin{proof}
	Užtenka pratikrinti kad pasirinkti algebros generatoriai $\vb e_{23}$ ir $\vb e_{31}$
	tapusavyje tenkina tokius pat sandaugos sąryšius kaip jų įvazdzdis.
	\begin{align}
		{\vb e_{23}}^2 = {\vb e_{31}}^2 &= - \mathbb1
		\\
		\vb e_{23} \vb e_{31} &= - \vb e_{12}
		\\
		\vb e_{31} \vb e_{23} &= \vb e_{12}
	\end{align}
	\begin{align}
		\qty(i\sigma_1)^2 = \qty(i\sigma_2)^2 = - {\sigma_2}^2 &= - I
		\\
		i\sigma_1 i\sigma_2 = - \sigma_1 \sigma_2 &= - i\sigma_3
		\\
		i\sigma_2 i\sigma_1 = - \sigma_2 \sigma_1 = \sigma_1 \sigma_2 &= i\sigma_3
	\end{align}
\end{proof}


\begin{preposition}
	\label{prep:pauli_rep-}
	Veikimu ant $\Clm$ \emph{lyginio} poalgebrio $A_0^-$ bazės apibrėžtas atvaizdavimas
	\begin{align*}
		\mathbb 1	\mapsto	I		\qquad
		\vb e_{23} 	\mapsto	-i\sigma_1	\qquad
		\vb e_{31}	\mapsto	-i\sigma_2	\qquad
		\vb e_{12}	\mapsto	-i\sigma_3	
	\end{align*}
	yra algebros $\Clm$ \emph{lyginio} poalgebrio ištikimas matricinis įvazdis.
\end{preposition}
\begin{proof}
	Analogiškai kaip teiginyje \ref{prep:pauli_rep+}
\end{proof}

\begin{preposition} 
	Atvizdis iš teiginio \ref{prep:pauli_rep+} siunčia $\spin$ į matricinė grupė $\mathup{SU}(2)$.
	Tai yra $\spin$ ir $\mathup{SU}(2)$ yra izomorfiškos grupės.
\end{preposition}
\begin{proof}
	Bendros formos $\spin$ elementas
	$ a_0 \mathbb1 + a_1 \vb e_{23} + a_2 \vb e_{31} + a_3 \vb e_{12}$,
	$ {a_0}^2 +{a_1}^2 +{a_2}^2 +{a_3}^2 = 1 $,
	yra atvaizduojamas į maricą
	\begin{equation}
		a_0 I + a_1 i\sigma_1 + a_2 i\sigma_2 + a_3 i\sigma_3
		=
		\mqty(
			  a_0 + ia_3	&& a_2 + ia_1 \\
			- a_2 + ia_1	&& a_0 - ia_3 
		)
		=
		\mqty(
			z_1	&& -z_2^* \\
			z_2	&& z_1^* 
		)
		\,,
	\end{equation}
	čia pasižymėjome kompleksinius skaičius 
	$z_1 = a_0 + i a_3$ ir $z_2 = - a_2 + ia_1$. 
	Sąlyga $ {a_0}^2 +{a_1}^2 +{a_2}^2 +{a_3}^2 = 1 $
	tada tampa $ \abs{z_1}^2 + \abs{z_2}^2 = 1$.
	Kas reiškia kad $
		\smqty(
			z_1	&& -z_2^* \\
			z_2	&& z_1^* 
		)
	$ yra benda forma matricos, kuri priklauso $\mathup{SU}(2)$.
\end{proof}

Gerai žinoma jog $\mathup{SU}(2)$ elmentas atitinkantis pasukimą
kampu $\theta$ aplink vienetinį vektoriu $\vb n$ yra
\begin{equation}
	e^{\frac{i\theta}{2} (\vb n,\symbf\sigma)}
	\,,
\end{equation}
čia $(\vb n,\symbf\sigma) = n_1 \sigma_1 + n_2 \sigma_2 + n_3 \sigma_3 $
ir kvadratinės matricos eksponetė yra apibrėžta kaip 
\begin{equation}
	e^X := \sum_{n=0}^\infty \frac{X^n}{n!}
	\,.
\end{equation}

Algebrų $\Clpm$ įvaizdžiams užrašyti pasinaudosime gama maricomis,
kurios chiralėje bazėja yra apibrežiamos kaip
\begin{equation}
	\gamma^0 := \mqty(\dmat{I,I}) 			\,,\quad
	\gamma^1 := \mqty(\admat{\sigma_1,-\sigma_1})	\,,\quad
	\gamma^2 := \mqty(\admat{\sigma_2,-\sigma_2})	\,,\quad
	\gamma^3 := \mqty(\admat{\sigma_3,-\sigma_3})	\,.
\end{equation}
Apsiribosime tiktai maricomis $\gamma^1$, $\gamma^2$ ir $\gamma^3$.
Jos tenkina sąryšius
\begin{equation}
		{\gamma^i}^2 = -\mqty(\dmat{I,I})
		\,,\quad
		\gamma^i \gamma^j = - \gamma^j \gamma^i
		\,.
\end{equation}


\begin{preposition}
	Veikimu ant $\Clm$ bazės apibrėžtas atvaizdavimas
	\begin{align*}
		\mathbb 1	\mapsto				&\mqty(\dmat{I,I})			&\quad	\vb e_{123}	\mapsto	\gamma^1\gamma^2\gamma^3 =	&\mqty(\admat{-iI,iI})			\\
		\vb e_{23}	\mapsto	\gamma^2 \gamma^3 =	&\mqty(\dmat{-i\sigma_1,-i\sigma_1})	&\quad	\vb e_1 	\mapsto	\gamma^1 =			&\mqty(\admat{\sigma_1,-\sigma_1})	\\
		\vb e_{31}	\mapsto	\gamma^3 \gamma^1 =	&\mqty(\dmat{-i\sigma_2,-i\sigma_2})	&\quad	\vb e_2 	\mapsto	\gamma^2 =			&\mqty(\admat{\sigma_2,-\sigma_2})	\\
		\vb e_{12}	\mapsto	\gamma^1 \gamma^2 =	&\mqty(\dmat{-i\sigma_3,-i\sigma_3})	&\quad	\vb e_3 	\mapsto	\gamma^3 =			&\mqty(\admat{\sigma_3,-\sigma_3})	\\
	\end{align*}
	yra algebros $\Clm$ ištikimas matricinis įvazdis.
\end{preposition}
\begin{proof}
	Gama matricos tenkina tokius pat daugybos sąryšius,
	kaip algebros $\Clm$ generatoriai $\vb e_1$, $\vb e_2$ ir $\vb e_3$.
	\begin{align}
		{\vb e_i}^2 &= -\mathbb1		&\quad	{\gamma^i}^2 &= - \mqty(\dmat{I,I})\\
		\vb e_i \vb e_j &= -\vb e_j \vb e_i	&\quad	\gamma^i\gamma^j &= - \gamma^j\gamma^i
	\end{align}
	Nesunku pamatyti, kad bazinių algebros vectorių atvaizdžiai
	yra tiesiškai neprikausomos matricos.
	Iš to seką kad $\gamma^1$, $\gamma^2$, $\gamma^3$
	ir $\vb e_1$, $\vb e_2$, $\vb e_3$
	generuoja izomorfiškas algebras,
	tai yra turime ištikimą matricinį įvaizdį.
\end{proof}

Verta pabrėžti, kad lyginiai bazės elementai atvaizduojami į blok-diagonalias matricas,
kur kiekvienas blokas yra kopija įvaizdžio iš teiginio \ref{prep:pauli_rep-}.
Tai lyginį $\pinm$ elementą, kuris atinitnka sukimą apie $\vb n$ kampu $\theta$,
gauname kaip
\begin{equation}
		\mqty(\dmat{
			e^{-\frac{i\theta}{2}(\vb n,\symbf\sigma)},
			e^{-\frac{i\theta}{2}(\vb n,\symbf\sigma)}
		})
		=
		\exp(\frac{\theta}{2}(
			 n_1\gamma^2\gamma^3
			+n_2\gamma^3\gamma^1
			+n_3\gamma^1\gamma^2
		))
		\,.
\end{equation}
Visi nelyginiai elementai gali būti gauti,
kaip lyginis elementas padaugintas iš $\gamma^1\gamma^2\gamma^3 = \smqty(&-iI\\iI&)$,
tai nelyginio elemeto bendra forma yra
\begin{equation}
		\mqty(\admat{
			-ie^{-\frac{i\theta}{2}(\vb n, \symbf\sigma)},
			 ie^{-\frac{i\theta}{2}(\vb n, \symbf\sigma)}
		})
		=
		\gamma^1\gamma^2\gamma^3 \exp(\frac{\theta}{2}(
			 n_1\gamma^2\gamma^3
			+n_2\gamma^3\gamma^1
			+n_3\gamma^1\gamma^2
		))
		\,.
\end{equation}


\begin{preposition}
	Veikimu ant $\Clp$ bazės apibrėžtas atvaizdavimas
	\begin{align*}
		\mathbb 1	\mapsto				&\mqty(\dmat{I,I})			&\quad	\vb e_{123}	\mapsto	-i\gamma^1\gamma^2\gamma^3 =	&\mqty(\admat{I,-I})			\\
		\vb e_{23}	\mapsto	-\gamma^2 \gamma^3 =	&\mqty(\dmat{i\sigma_1,i\sigma_1})	&\quad	\vb e_1 	\mapsto	i\gamma^1 =			&\mqty(\admat{i\sigma_1,-i\sigma_1})	\\
		\vb e_{31}	\mapsto	-\gamma^3 \gamma^1 =	&\mqty(\dmat{i\sigma_2,i\sigma_2})	&\quad	\vb e_2 	\mapsto	i\gamma^2 =			&\mqty(\admat{i\sigma_2,-i\sigma_2})	\\
		\vb e_{12}	\mapsto	-\gamma^1 \gamma^2 =	&\mqty(\dmat{i\sigma_3,i\sigma_3})	&\quad	\vb e_3 	\mapsto	i\gamma^3 =			&\mqty(\admat{i\sigma_3,-i\sigma_3})	\\
	\end{align*}
	yra algebros $\Clp$ ištikimas matricinis įvazdis.
\end{preposition}
\begin{proof}
	Patikriname ar įvaizdis gerbia generatorių $\vb e_1$, $\vb e_2$ ir $\vb e_3$ 
	daugybos sąryšius.
	\begin{align}
		{\vb e_i}^2 &= \mathbb1			&\quad	{\gamma^i}^2 &= \mqty(\dmat{I,I})\\
		\vb e_i \vb e_j &= -\vb e_j \vb e_i	&\quad	(i\gamma^i)(i\gamma^j) &= - (i\gamma^j)(i\gamma^i)
	\end{align}.
	Matome, kad bazė atvaizduojama į tieškai neprikausomas matricas,
	tai turime ištikimą matricinį algebros $\Clp$ įvaizdį.
\end{proof}

Šį kartą lyginiai bazės elementai yra atviazduojami į blok-diagonamias matricas,
kurios yra dvi kopijos įvaiždžio iš teiginio \ref{prep:pauli_rep+} ir lyginio elemento bendra forma yra
\begin{equation}
		\mqty(\dmat{
			e^{\frac{i\theta}{2}(\vb n , \symbf\sigma)},
			e^{\frac{i\theta}{2}(\vb n , \symbf\sigma)}
		})
		=
		\exp(-\frac{\theta}{2}(
			 n_1\gamma^2\gamma^3
			+n_2\gamma^3\gamma^1
			+n_3\gamma^1\gamma^2
		))
		\,.
\end{equation}
Nelyginio elemeto bendra forma yra
\begin{equation}
		\mqty(\admat{
			 e^{\frac{i\theta}{2}(\vb n , \symbf\sigma)},
			-e^{\frac{i\theta}{2}(\vb n , \symbf\sigma)}
		})
		=
		-\gamma^1\gamma^2\gamma^3 \exp(-\frac{\theta}{2}(
			 n_1\gamma^2\gamma^3
			+n_2\gamma^3\gamma^1
			+n_3\gamma^1\gamma^2
		))
		\,.
\end{equation}






